% mazzi: 04
\chapter{Le leggi regionali in materia di impiantistica sportiva}

\section{Il rinvio alla legge regionale}

\textbf{Art.\ 90 comma 25 L.\ 289/2002}: quando l'ente pubblico territoriale --- nella maggior
parte dei casi il \textbf{Comune} proprietario dell'impianto --- non intenda gestire direttamente
gli impianti sportivi, la gestione è affidata \textbf{in via preferenziale} a società e
associazioni sportive dilettantistiche, sulla base di \textbf{convenzioni} che ne stabiliscono i
criteri d'uso e previa determinazione di criteri generali e obiettivi per l'individuazione dei
soggetti affidatari; le \textbf{Regioni disciplinano, con propria legge, le modalità di
affidamento}.

\begin{itemize}
  \item è una \textbf{norma quadro} proprio perché rinvia a una normativa regionale di disciplina
    delle modalità di affidamento: la linea guida statale è la preferenzialità, la sua
    declinazione concreta spetta alle Regioni;
  \item introduce il principio del \textbf{convenzionamento}: la convenzione è un atto \textit{iure
    privatorum}, una scrittura privata, che stabilisce i criteri d'uso dell'impianto;
  \item i criteri generali e obiettivi per l'individuazione degli affidatari vanno
    \textbf{predeterminati}, a monte della procedura.
\end{itemize}

\rubrica{Le tredici leggi regionali emanate}

\begin{multicols}{2}
\begin{enumerate}
  \item L.R.\ Valle d'Aosta 4 agosto 2006, n.\ 18;
  \item L.R.\ Puglia 4 dicembre 2006, n.\ 33;
  \item L.R.\ Lombardia 14 dicembre 2006, n.\ 27;
  \item L.R.\ Umbria 12 marzo 2007, n.\ 5;
  \item L.R.\ Liguria 7 ottobre 2009, n.\ 40;
  \item L.R.\ Calabria 22 novembre 2010, n.\ 28;
  \item L.R.\ Molise 9 settembre 2011, n.\ 18;
  \item L.R.\ Marche 2 aprile 2012, n.\ 5;
  \item L.R.\ Abruzzo 19 giugno 2012, n.\ 27;
  \item L.R.\ Campania 25 novembre 2013, n.\ 18;
  \item L.R.\ Toscana 6 marzo 2015, n.\ 21;
  \item L.R.\ Veneto 15 maggio 2015, n.\ 8;
  \item L.R.\ Emilia-Romagna 31 maggio 2017, n.\ 8.
\end{enumerate}
\end{multicols}

% ---------------------------------------------------------------------------

\section{La Valle d'Aosta: L.R.\ 18/2006, art.\ 10}

Regione a statuto speciale, la prima a legiferare: la legge è percorsa da passaggi che tendono a
\textbf{privilegiare le società del territorio}.

\rubrica{I criteri di affidamento (comma 2)}

Gli enti locali affidano gli impianti sportivi in gestione nel rispetto dei seguenti criteri:

\begin{enumerate}
  \item[a)] \textbf{apertura dell'impianto a tutti i cittadini} e \textbf{imparzialità} nel
    permetterne l'utilizzo, nel rispetto della \textbf{specificità di ogni impianto} e
    compatibilmente con l'attività dei soggetti affidatari;
  \item[b)] \textbf{avviso pubblico} come modalità di pubblicità nella procedura di selezione,
    idonea a garantirne l'effettiva conoscenza a tutti i soggetti interessati;
  \item[c)] \textbf{scelta dell'affidatario} che tenga conto di esperienza nel settore,
    radicamento sul territorio nel bacino di utenza dell'impianto, affidabilità economica,
    qualificazione professionale di istruttori e allenatori impiegati, compatibilità dell'attività
    sportiva esercitata con quella praticabile nell'impianto ed eventuale organizzazione di
    attività a favore di giovani, disabili e anziani;
  \item[d)] \textbf{selezione} da effettuare valutando sia i profili gestionali, tecnici e sociali
    dell'offerta, sia la sua convenienza economica, sulla base della previa indicazione, da parte
    dell'ente locale, del \textbf{canone minimo} che intende percepire e dell'eventuale
    \textbf{contributo massimo} che intende concedere a sostegno della gestione;
  \item[e)] \textbf{compatibilità} delle eventuali attività ricreative e sociali di interesse
    pubblico, praticabili straordinariamente negli impianti, con il normale uso sportivo;
  \item[f)] determinazione della \textbf{durata massima} dell'affidamento in gestione.
\end{enumerate}

Alcune precisazioni sui criteri:

\begin{itemize}
  \item l'apertura a tutti i cittadini riprende l'\textbf{art.\ 90 comma 24} L.\ 289/2002: gli
    impianti sportivi pubblici vanno messi e mantenuti a disposizione di un'utenza pubblica;
  \item la \textbf{specificità dell'impianto} restringe di fatto la platea dei potenziali
    affidatari alle società che praticano quella disciplina: vale per gli impianti non
    polisportivi, come gli impianti per gli sport del ghiaccio, un impianto natatorio o una pista
    di atletica;
  \item giovani, disabili e anziani sono trattati come \textbf{categorie privilegiate} dalla norma;
  \item l'affidabilità economica fa emergere il profilo della \textbf{sostenibilità economico-%
    finanziaria} della gestione;
  \item \textbf{canone} e \textbf{contributo} sono le due grandezze speculari del rapporto: il
    canone minimo è corrisposto dal soggetto affidatario all'ente, il contributo massimo è quanto
    l'amministrazione può elargire per garantire l'\textbf{equilibrio economico-finanziario} della
    gestione.
\end{itemize}

\rubrica{La convenzione (comma 3)}

Gli enti locali stipulano con il soggetto affidatario una convenzione concernente la gestione
dell'impianto, che stabilisce in particolare i \textbf{criteri d'uso} dell'impianto e i criteri
\textbf{anche economici} della gestione.

\rubrica{La deroga per gli impianti destinati a specifiche discipline (comma 4)}

In deroga al comma 2 --- e quindi all'avviso pubblico --- gli enti locali, \textbf{sentito il
CONI}, hanno \textbf{facoltà} (non obbligo) di affidare gli impianti sportivi \textbf{privi di
rilevanza economica} e \textbf{destinati a specifiche discipline}, stipulando apposita convenzione:

\begin{itemize}
  \item con la \textbf{federazione sportiva nazionale corrispondente};
  \item ovvero con una \textbf{società o associazione sportiva dilettantistica sua affiliata}, da
    essa indicata, della quale la Federazione o il CONI si faccia \textbf{garante} degli obblighi
    previsti in convenzione nei confronti dell'ente locale;
\end{itemize}

resta fermo in ogni caso il comma 2 lettera a), cioè l'apertura dell'impianto a tutti i cittadini.

L'esempio classico è lo \textbf{stadio del ghiaccio}; non lo è il campo da calcio, sulle cui
misure sono spesso praticabili anche altre discipline, il che renderebbe difficoltoso
l'affidamento a una singola federazione. È previsione da riconsiderare alla luce dell'evoluzione
normativa successiva al 2006, della prassi amministrativa, delle delibere ANAC e della
giurisprudenza.

% ---------------------------------------------------------------------------

\section{L'affidamento diretto alle Federazioni sportive nazionali}

Il meccanismo introdotto dalla Valle d'Aosta è stato poi utilizzato anche altrove, in altri
contesti, come \textbf{escamotage} per evitare la gara. I limiti di ammissibilità:

\begin{itemize}
  \item è ammissibile \textbf{solo} se l'impianto \textbf{non è idoneo} alla pratica di discipline
    regolamentate da altre FSN --- cioè non è omologabile per altre attività agonistiche --- e
    \textbf{è privo di rilevanza economica};
  \item consente di evitare il ricorso alla \textbf{gara informale} mediante invito di minimo
    \textbf{5 concorrenti} (previsione del vecchio codice dei contratti), ma non può in alcun modo
    derogare ai \textbf{principi comunitari di trasparenza e non discriminazione};
  \item va valutata la presenza, sul territorio di riferimento, di \textbf{enti di promozione
    sportiva} (EPS) o loro affiliati che pratichino la disciplina cui l'impianto è destinato e che
    non siano affiliati alla federazione di riferimento: la questione non si esaurisce quindi
    guardando alla sola FSN;
  \item ove la FSN \textbf{non gestisse direttamente} l'impianto, sarebbe parimenti tenuta a
    esperire una \textbf{procedura ad evidenza pubblica} per l'affidamento in gestione.
\end{itemize}

\rubrica{Il fondamento: art.\ 23 comma 1 Statuto CONI}

``\dots hanno \textbf{valenza pubblicistica} le attività delle Federazioni sportive nazionali
relative \dots all'utilizzazione e alla gestione degli impianti sportivi pubblici\dots''

\begin{itemize}
  \item le FSN sono \textbf{associazioni di diritto privato}, ma a questi fini vanno considerate
    \textbf{organismi di diritto pubblico}: lo Statuto CONI elenca una serie di attività delegate
    dal CONI, ente pubblico, alle federazioni, che hanno appunto valenza pubblicistica;
  \item in questi ambiti le federazioni devono ottemperare ai criteri e ai principi di trasparenza
    della \textbf{L.\ 241/1990} sul procedimento amministrativo e agli adempimenti conseguenti;
  \item ne segue che la normativa valdostana è \textbf{imprudente}, o quantomeno audace, nel
    prevedere l'affidamento diretto a federazioni o ad affiliati da esse indicati senza chiarire
    che anche le federazioni sono tenute a una \textbf{procedura comparativa} fra i propri
    affiliati, o quantomeno fra quelli del territorio di riferimento.
\end{itemize}

% ---------------------------------------------------------------------------

\section{Dalla Puglia all'Emilia-Romagna}

\subsection{Puglia --- L.R.\ 33/2006, art.\ 18 (mod.\ L.R.\ 32/2012)}

Attua l'art.\ 90 comma 25 disciplinando le modalità di affidamento a terzi degli impianti sportivi
di proprietà degli enti pubblici territoriali.

\begin{itemize}
  \item \textbf{ambito di applicazione}: gli impianti di proprietà di enti pubblici territoriali
    \textbf{non gestiti direttamente} dagli enti medesimi, intesi quali strutture in cui possono
    praticarsi attività sportive \textbf{di qualsiasi livello}, eventualmente associate ad attività
    ricreative e sociali di interesse pubblico;
  \item ribadisce il principio del comma 24: l'uso degli impianti è \textbf{aperto e accessibile a
    tutte le cittadine e a tutti i cittadini};
  \item i soggetti affidatari sono individuati con \textbf{procedure di evidenza pubblica} fra
    coloro che presentano idonei requisiti e garantiscono le finalità dell'art.\ 18;
  \item la gestione è affidata \textbf{in via preferenziale}, \textbf{favorendone l'aggregazione},
    a FSN, discipline sportive associate (DSA), EPS riconosciuti dal CONI, società e associazioni
    sportive dilettantistiche con i requisiti dell'art.\ 90 L.\ 289/2002 e successivi regolamenti
    attuativi --- quindi iscritte al registro CONI;
  \item l'uso dell'impianto va garantito \textbf{anche ai soggetti non affidatari}, purché in
    possesso degli stessi requisiti: vale cioè il principio dell'utilizzo plurimo.
\end{itemize}

La locuzione ``in via preferenziale'' è riportata \textit{sic et simpliciter} dalla norma
nazionale, \textbf{senza declinare} il concetto di preferenzialità: sotto il profilo
dell'affidamento la legge pugliese non offre spunti di novità.

\subsection{Lombardia --- L.R.\ 27/2006}

È la legge che più di ogni altra \textbf{differenzia la procedura di selezione} a seconda che si
tratti di impianto con o senza rilevanza economica. Vi ha inciso il dibattito in corso all'epoca
dell'approvazione sull'affidamento dei servizi privi di rilevanza economica, poi sfociato in
modifiche normative e in un quesito referendario.

\begin{itemize}
  \item \textbf{impianti con rilevanza economica} --- cioè la cui gestione è capace di produrre
    reddito --- che per dimensioni e altre caratteristiche richiedono una \textbf{gestione di tipo
    imprenditoriale}: le società e associazioni sportive dilettantistiche devono dimostrare di
    avere i \textbf{requisiti imprenditoriali e tecnici necessari}. La qualifica di ASD o SSD non è
    di per sé sufficiente: la preferenzialità resta, ma condizionata alla capacità di una gestione
    manageriale, spesso richiesta da impianti a tecnologia complessa;
  \item \textbf{impianti senza rilevanza economica}: \textbf{affidamento diretto} ammesso anche
    \textbf{agli utilizzatori} degli impianti stessi.
\end{itemize}

\subsection{Umbria --- L.R.\ 5/2007, art.\ 4}

\rubrica{Affidamento diretto per gli impianti privi di rilevanza economica (comma 4)}

Ammesso nei casi seguenti:

\begin{enumerate}
  \item[a)] quando sul territorio di riferimento è presente un \textbf{solo soggetto} che promuove
    la disciplina sportiva praticabile presso l'impianto --- una gara sarebbe inutile;
  \item[b)] quando il servizio è affidato direttamente a una \textbf{società a capitale interamente
    o a maggioranza pubblica}, costituita nelle forme dell'art.\ 90 comma 17 lettera c) L.\
    289/2002. Se il capitale è solo a maggioranza pubblica, e non interamente pubblico, si
    applicano le norme sulle \textbf{società \textit{in house}} e occorre comunque una procedura
    selettiva per individuare il socio privato;
  \item[c)] quando i soggetti sportivi operanti sul territorio su cui insiste l'impianto
    costituiscono un \textbf{unico soggetto sportivo}: raggruppamento temporaneo, polisportiva di
    secondo livello o altra forma di aggregazione fra tutti i soggetti che avrebbero i requisiti
    per ambire alla gestione. Anche qui la gara può essere bypassata, ma l'affidamento diretto
    dev'essere preceduto da forme di pubblicità sulla volontà dell'amministrazione di aggiudicare
    l'impianto.
\end{enumerate}

\rubrica{Impianti con rilevanza economica (commi 5--6)}

\begin{itemize}
  \item l'affidamento avviene con \textbf{procedura ad evidenza pubblica}, aperta anche ai soggetti
    dell'art.\ 90 comma 17 L.\ 289/2002 e ai loro raggruppamenti, nel rispetto dei criteri del
    comma 2 lettere a), b), c) e dei vincoli statali e comunitari;
  \item sono ammesse a partecipare anche \textbf{società di capitali}, pure in aggregazione con i
    soggetti dell'art.\ 90 comma 17 lettere a) e b): soggetti \textit{profit} che concorrono
    insieme ad associazioni e società sportive dilettantistiche.
\end{itemize}

\subsection{Liguria --- L.R.\ 40/2009, art.\ 19}

Il tratto caratterizzante è un'\textbf{esclusione} dall'applicazione del capo:

\begin{enumerate}
  \item[a)] gli impianti situati in \textbf{comuni con popolazione inferiore ai 5.000 abitanti}, a
    eccezione delle piscine e delle sale con caratteristiche di palazzi dello sport;
  \item[b)] gli impianti facenti parte del \textbf{patrimonio regionale}, affidati in gestione
    direttamente a enti pubblici.
\end{enumerate}

Resta garantito a tutti i cittadini l'uso degli impianti sportivi pubblici (art.\ 2).

La prima esclusione è tecnicamente e giuridicamente discutibile, pur comprendendone la
\textit{ratio}: un impianto a rilevanza economica può ben trovarsi nel territorio di un comune
sotto i 5.000 abitanti pur ricadendo in un'area metropolitana, e l'eccezione limitata a piscine e
palazzi dello sport non è sufficiente a coprire l'ipotesi.

\subsection{Sardegna --- L.R.\ 17/1999, art.\ 21}

Regione autonoma, con una previsione di legge e un regolamento attuativo. È legge
\textbf{anteriore} all'art.\ 90 L.\ 289/2002, in vigore dal 2003, e la sua datazione traspare
dalla terminologia usata.

\rubrica{Regolamento di gestione degli impianti}

\begin{itemize}
  \item gli enti e le associazioni beneficiarie delle \textbf{provvidenze} previste dalla legge
    --- i contributi --- sono tenuti ad adottare un \textbf{regolamento di gestione} degli
    impianti, formulato sulla base del \textbf{regolamento tipo} predisposto dall'Assessorato
    regionale competente in materia di sport e approvato dalla Giunta regionale;
  \item nei limiti della fruibilità degli impianti e delle relative attrezzature, l'uso va
    garantito ai \textbf{sodalizi sportivi} e alle associazioni del tempo libero operanti nel
    territorio interessato, nonché alla \textbf{popolazione scolastica} che non disponga di
    adeguate strutture;
  \item le stesse disposizioni si applicano alle opere realizzate in base alla preesistente
    normativa regionale in materia di sport.
\end{itemize}

\rubrica{Criteri per l'assegnazione (regolamento ex art.\ 21)}

La \textbf{priorità} nella scelta del concessionario è data agli operatori sportivi che già
svolgono attività nella disciplina praticata nell'impianto e nell'ambito del territorio della
circoscrizione interessata, tenendo conto in via prioritaria di:

\begin{multicols}{2}
\begin{itemize}
  \item numero degli atleti tesserati;
  \item anni di attività del sodalizio;
  \item livello dei campionati cui partecipa;
  \item risultati agonistici ottenuti;
  \item attività di promozione dello sport fra i giovani in età scolare.
\end{itemize}
\end{multicols}

La norma parla di \textbf{priorità} e non di preferenzialità --- è anteriore all'art.\ 90 --- ma il
concetto è il medesimo. La previsione sconta però la disciplina dell'epoca e non è più adeguata
agli standard del codice dei contratti pubblici né all'orientamento di derivazione comunitaria,
anche in materia di concessioni demaniali: non è più possibile riconoscere un \textbf{diritto di
insistenza} a chi già gestisce l'impianto o vi opera.

\subsection{Molise --- L.R.\ 18/2011, art.\ 2}

\begin{itemize}
  \item gli enti pubblici territoriali che non intendano gestire direttamente i propri impianti,
    anche mediante convenzione fra gli enti stessi, ne affidano \textbf{in via preferenziale} la
    gestione a società sportive, ASD, EPS, DSA e FSN, anche in forma associata;
  \item l'affidamento avviene \textbf{prioritariamente} a vantaggio dei soggetti che praticano la
    disciplina cui l'impianto è destinato e che operano nel territorio dell'ente affidatario,
    secondo il \textbf{criterio territoriale}: comune, provincia, regione;
  \item \textbf{elemento di novità}: l'affidamento a \textbf{soggetti diversi} avviene solo in caso
    di \textbf{esito infruttuoso} delle procedure di selezione dell'art.\ 3, e comunque nel
    rispetto dei relativi principi --- prima una procedura fra i soggetti sportivi indicati da
    legge statale e regionale, poi, solo se andata deserta, l'apertura ad altri;
  \item per gli \textbf{impianti con rilevanza economica} che richiedono una gestione di tipo
    imprenditoriale l'affidamento avviene con \textbf{procedura di evidenza pubblica} e i soggetti
    sportivi devono dimostrare di possedere i \textbf{requisiti economici e tecnici} necessari ---
    come già previsto dalla legge lombarda.
\end{itemize}

\subsection{Marche --- L.R.\ 5/2012, art.\ 19}

\begin{itemize}
  \item i soggetti affidatari sono individuati fra coloro che presentano idonei requisiti, con
    \textbf{procedure di evidenza pubblica} nel rispetto della normativa vigente;
  \item la gestione è affidata \textbf{in via preferenziale} a SSD e ASD, EPS, DSA e FSN;
  \item gli enti territoriali stipulano con gli affidatari \textbf{convenzioni} che stabiliscono i
    criteri d'uso degli impianti;
  \item l'uso dell'impianto è garantito \textbf{anche alle società e associazioni non
    affidatarie}.
\end{itemize}

\subsection{Campania --- L.R.\ 18/2013, art.\ 20}

Gli enti pubblici territoriali che non gestiscono direttamente, nel rispetto del principio di
\textbf{imparzialità della scelta}, affidano la gestione al CONI, al \textbf{CIP} --- Comitato
italiano paralimpico, ente pubblico esplicitamente menzionato accanto al CONI --- alle FSN, agli
EPS o alle DSA, alle ASD iscritte alla sezione A prevista dall'art.\ 11 comma 2, che garantiscono:

\begin{enumerate}
  \item[a)] l'\textbf{apertura dell'impianto} a tutti i soggetti e, per le piscine, adeguati spazi
    per il \textbf{nuoto libero};
  \item[b)] esperienza nel settore, radicamento nel territorio del bacino di utenza, affidabilità
    economica, qualificazione professionale di istruttori e operatori;
  \item[c)] compatibilità dell'attività sportiva esercitata con quella praticabile nell'impianto e
    con l'organizzazione di attività a favore di giovani, diversamente abili e anziani;
  \item[d)] compatibilità delle attività ricreative e sociali d'interesse pubblico con il normale
    uso sportivo;
  \item[e)] lo svolgimento dell'\textbf{attività agonistica};
  \item[f)] la destinazione di \textbf{investimenti} alla migliore fruizione dell'impianto.
\end{enumerate}

\rubrica{Procedura (commi 3--4)}

\begin{itemize}
  \item \textbf{avviso pubblico} come modalità di pubblicità della selezione, trasmesso alla
    competente struttura regionale che ne cura la pubblicazione nel \textbf{Bollettino Ufficiale
    della Regione};
  \item la selezione avviene in ragione della diversa tipologia di impianto, nel rispetto dei
    criteri di \textbf{proporzionalità, non discriminazione e trasparenza}, all'esito della
    valutazione dei \textbf{progetti} presentati, che consentono di apprezzare i profili economici
    e tecnici della gestione;
  \item si tiene conto anche del \textbf{canone minimo} che l'ente intende percepire e del
    \textbf{contributo massimo} che intende concedere a sostegno della gestione.
\end{itemize}

\rubrica{Affidamento diretto nei comuni sotto i 5.000 abitanti}

In deroga, il servizio può essere affidato in via diretta dai comuni con popolazione inferiore ai
cinquemila abitanti:

\begin{enumerate}
  \item[a)] quando gli impianti hanno caratteristiche e dimensioni che consentono lo svolgimento di
    attività \textbf{esclusivamente amatoriali e ricreative} --- non omologabili per l'attività
    agonistica --- riferibili al territorio dove sono ubicati;
  \item[b)] quando nel territorio di riferimento dell'ente proprietario è presente un \textbf{solo
    soggetto} che promuove la disciplina praticabile presso l'impianto;
  \item[c)] quando le società e le associazioni di promozione sportiva operanti nel territorio
    costituiscono un \textbf{unico soggetto sportivo}, anche in forma associativa o consortile.
\end{enumerate}

\subsection{Toscana --- L.R.\ 21/2015, art.\ 14}

\begin{itemize}
  \item gli enti locali che non intendono gestire direttamente affidano la gestione \textbf{in via
    preferenziale} a ASD e SSD, EPS, DSA e FSN, secondo \textbf{procedure ad evidenza pubblica};
  \item l'affidamento a \textbf{soggetti diversi} avviene solo in caso di \textbf{esito
    infruttuoso} delle procedure di selezione previste, comunque nel rispetto dei relativi
    principi --- come già nel Molise.
\end{itemize}

\subsection{Veneto --- L.R.\ 8/2015, artt.\ 25--26}

\begin{itemize}
  \item affidamento \textbf{in via preferenziale} a società e associazioni sportive
    dilettantistiche \textbf{senza fini di lucro}, FSN, DSA ed EPS, nonché a \textbf{consorzi e
    associazioni} fra i predetti soggetti;
  \item affidamento a soggetti diversi, in possesso di idonei requisiti, \textbf{esclusivamente}
    in caso di esito infruttuoso delle modalità dell'art.\ 26;
  \item l'individuazione avviene con \textbf{procedure ad evidenza pubblica} nel rispetto dei
    principi di imparzialità, trasparenza e adeguata pubblicità;
  \item \textbf{affidamento in via diretta} ammesso quando ricorre almeno uno dei presupposti:
    \begin{enumerate}
      \item[a)] presenza sul territorio di riferimento dell'ente locale di un \textbf{solo
        soggetto} che promuova la disciplina praticabile nell'impianto;
      \item[b)] presenza sul territorio di società e altri soggetti di promozione sportiva
        operanti tramite un \textbf{unico soggetto sportivo}.
    \end{enumerate}
\end{itemize}

\subsection{Emilia-Romagna --- L.R.\ 8/2017}

È l'ultima legge regionale approvata in materia. Gli enti locali disciplinano condizioni e
modalità di affidamento in gestione degli impianti, con particolare riferimento a quelli aventi
\textbf{minore rilevanza economica} --- formula che si colloca fra le due categorie consuete e
allude agli impianti a redditività limitata --- sulla base dei principi seguenti:

\begin{enumerate}
  \item[a)] \textbf{compatibilità} fra le attività sportive praticabili e quelle esercitate negli
    impianti, favorendone l'uso da parte dei praticanti del territorio;
  \item[b)] \textbf{valorizzazione delle potenzialità} degli impianti, con un rapporto equilibrato
    --- ove compatibile con le caratteristiche degli impianti --- fra normale uso sportivo,
    utilizzazione da parte del pubblico, attività di promozione della pratica sportiva e attività
    ricreative e sociali;
  \item[c)] valutazione dei requisiti di \textbf{qualificazione e affidabilità economica}, nonché
    delle competenze e capacità maturate in precedenti esperienze di gestione;
  \item[d)] valutazione dell'offerta sulla base del \textbf{miglior rapporto fra qualità e prezzo}
    --- l'offerta economicamente più vantaggiosa --- secondo criteri predeterminati, purché sia
    assicurato l'\textbf{equilibrio economico} della gestione.
\end{enumerate}

La Regione individua inoltre \textbf{linee guida} contenenti migliori pratiche \textbf{non
vincolanti}, ai fini della loro promozione sul territorio, e le correlate definizioni applicative.

% ---------------------------------------------------------------------------

\section{La competenza all'affidamento: Giunta o Consiglio comunale?}

\textbf{TAR Calabria 1 luglio 2010, n.\ 141}: risolve il dubbio muovendo dall'\textbf{art.\ 42
TUEL}, che definisce il \textbf{Consiglio comunale} ``l'organo di indirizzo e di controllo
politico-amministrativo'' dell'ente, attribuendogli ``la competenza in materia di appalti
(\dots) lasciando il concreto affidamento dell'appalto alla competenza della \textbf{Giunta}, in
quanto attività meramente esecutiva o di ordinaria amministrazione''.

% ---------------------------------------------------------------------------

\section{Sintesi della disciplina applicabile}

\rubrica{I criteri ricorrenti nelle leggi regionali}

\begin{itemize}
  \item \textbf{preferenzialità declinata come priorità procedurale}: Molise, Toscana e Veneto
    riservano una prima procedura ai soli soggetti sportivi e aprono ad altri solo in caso di esito
    infruttuoso. È il tentativo più esplicito di dare contenuto operativo alla locuzione dell'art.\
    90 comma 25, ed è già stato oggetto di \textbf{censura da parte della Commissione europea} per
    violazione delle norme e dei principi sulla concorrenza;
  \item \textbf{preferenzialità non declinata}: la Puglia riprende la formula statale senza
    precisarla;
  \item \textbf{differenziazione per rilevanza economica}: Lombardia, Umbria e Molise distinguono
    la procedura a seconda che l'impianto abbia o meno rilevanza economica, richiedendo requisiti
    imprenditoriali e tecnici quando la gestione dev'essere di tipo imprenditoriale;
  \item \textbf{presupposti tipici dell'affidamento diretto}: unico soggetto che promuove la
    disciplina sul territorio; aggregazione di tutti i soggetti sportivi in un unico soggetto;
    società a capitale pubblico (Umbria); impianto privo di rilevanza economica destinato a
    specifica disciplina (Valle d'Aosta);
  \item \textbf{canone minimo e contributo massimo} come parametri della valutazione economica
    dell'offerta (Valle d'Aosta, Campania);
  \item \textbf{comuni sotto i 5.000 abitanti}: la Liguria li \textbf{esclude} dall'applicazione
    della disciplina, la Campania consente loro l'\textbf{affidamento diretto} a condizioni
    determinate.
\end{itemize}

\rubrica{La conclusione}

Il criterio di preferenzialità dettato dall'art.\ 90 comma 25 L.\ 289/2002 si porrebbe in evidente
violazione:

\begin{itemize}
  \item dei \textbf{principi comunitari} di libera circolazione delle merci, libertà di
    stabilimento, libera prestazione di servizi e soprattutto trasparenza, parità di trattamento e
    non discriminazione;
  \item delle \textbf{nuove norme interne} in materia di contrattualistica pubblica;
\end{itemize}

ma solo \textbf{laddove si traducesse in un semplicistico affidamento diretto} senza ricorrere ad
alcuna forma di \textbf{pubblicità preventiva}, che è dunque necessaria in ogni caso.

Anche quando l'impianto \textbf{non} avesse rilevanza economica, l'affidamento non potrebbe che
avvenire nel rispetto dei principi desumibili dal \textbf{Trattato} --- art.\ 30 D.Lgs.\ 163/2006,
per come poi riportato anche nel nuovo codice dei contratti pubblici.
